% ============================================================
%  notas_desarrollo.tex
%  Diario de desarrollo — TFG: Optimizacion SLAM TurtleBot4
%  Autor: Jesús Ruslan Esquinas Herrera
%
%  USO:
%    Añade una \entrada{} cada vez que resuelvas algo relevante.
%    Cuando escribas la memoria final, busca la \secciondestino
%    de cada entrada para saber dónde encaja el contenido.
%
%  Para incluirlo en tu documento principal (opcional):
%    % ============================================================
%  notas_desarrollo.tex
%  Diario de desarrollo — TFG: Optimizacion SLAM TurtleBot4
%  Autor: Jesús Ruslan Esquinas Herrera
%
%  USO:
%    Añade una \entrada{} cada vez que resuelvas algo relevante.
%    Cuando escribas la memoria final, busca la \secciondestino
%    de cada entrada para saber dónde encaja el contenido.
%
%  Para incluirlo en tu documento principal (opcional):
%    % ============================================================
%  notas_desarrollo.tex
%  Diario de desarrollo — TFG: Optimizacion SLAM TurtleBot4
%  Autor: Jesús Ruslan Esquinas Herrera
%
%  USO:
%    Añade una \entrada{} cada vez que resuelvas algo relevante.
%    Cuando escribas la memoria final, busca la \secciondestino
%    de cada entrada para saber dónde encaja el contenido.
%
%  Para incluirlo en tu documento principal (opcional):
%    % ============================================================
%  notas_desarrollo.tex
%  Diario de desarrollo — TFG: Optimizacion SLAM TurtleBot4
%  Autor: Jesús Ruslan Esquinas Herrera
%
%  USO:
%    Añade una \entrada{} cada vez que resuelvas algo relevante.
%    Cuando escribas la memoria final, busca la \secciondestino
%    de cada entrada para saber dónde encaja el contenido.
%
%  Para incluirlo en tu documento principal (opcional):
%    \input{notas_desarrollo}
% ============================================================

\documentclass[12pt, a4paper]{article}

\usepackage[utf8]{inputenc}
\usepackage[T1]{fontenc}
\usepackage[spanish]{babel}
\usepackage{geometry}
\geometry{margin=2.5cm}

\usepackage{xcolor}
\usepackage{mdframed}
\usepackage{listings}
\usepackage{hyperref}
\usepackage{parskip}
\usepackage{titlesec}
\usepackage{enumitem}
\usepackage{fancyhdr}
\usepackage{booktabs}
\usepackage{array}

% ── Colores ──────────────────────────────────────────────────
\definecolor{colorProblema}{RGB}{180, 60, 40}
\definecolor{colorSolucion}{RGB}{30, 120, 60}
\definecolor{colorDestino}{RGB}{30, 80, 160}
\definecolor{colorNota}{RGB}{140, 90, 0}
\definecolor{colorFondo}{RGB}{248, 248, 248}
\definecolor{colorBorde}{RGB}{200, 200, 200}
\definecolor{codeGray}{RGB}{245, 245, 245}

% ── Estilo de código ─────────────────────────────────────────
\lstset{
  backgroundcolor=\color{codeGray},
  basicstyle=\ttfamily\small,
  breaklines=true,
  frame=single,
  rulecolor=\color{colorBorde},
  numbers=left,
  numberstyle=\tiny\color{gray},
  tabsize=2,
  showstringspaces=false,
  keywordstyle=\color{colorDestino}\bfseries,
  commentstyle=\color{colorSolucion}\itshape,
  stringstyle=\color{colorProblema},
}

% ── Entorno de entrada de diario ─────────────────────────────
%  \entrada{fecha}{título}
\newenvironment{entrada}[2]{%
  \begin{mdframed}[
    backgroundcolor=colorFondo,
    linecolor=colorBorde,
    linewidth=0.8pt,
    topline=true, bottomline=true,
    leftline=true, rightline=false,
    leftmargin=0pt, rightmargin=0pt,
    innerleftmargin=10pt, innerrightmargin=10pt,
    innertopmargin=8pt, innerbottommargin=8pt,
  ]
  \noindent
  {\small\color{gray}\texttt{#1}} \hfill
  {\small\color{gray}---}\\[2pt]
  {\large\bfseries #2}\\[6pt]
  \hrule height 0.4pt \color{colorBorde}
  \vspace{6pt}
}{%
  \end{mdframed}
  \vspace{8pt}
}

% ── Comandos de etiquetas ─────────────────────────────────────
\newcommand{\problema}[1]{%
  \vspace{4pt}
  \noindent{\color{colorProblema}\bfseries Problema:} #1\par
}

\newcommand{\causa}[1]{%
  \vspace{2pt}
  \noindent{\color{colorProblema}\bfseries Causa:} #1\par
}

\newcommand{\solucion}[1]{%
  \vspace{2pt}
  \noindent{\color{colorSolucion}\bfseries Solución:} #1\par
}

\newcommand{\secciondestino}[1]{%
  \vspace{4pt}
  \noindent{\color{colorDestino}\bfseries $\rightarrow$ Sección en memoria:} \texttt{#1}\par
}

\newcommand{\nota}[1]{%
  \vspace{4pt}
  \noindent{\color{colorNota}\bfseries Nota:} \textit{#1}\par
}

\newcommand{\resultado}[1]{%
  \vspace{4pt}
  \noindent{\color{colorSolucion}\bfseries Resultado medible:} #1\par
}

% ── Cabecera ─────────────────────────────────────────────────
\pagestyle{fancy}
\fancyhf{}
\rhead{\small TFG — Notas de desarrollo}
\lhead{\small Jesús Ruslan}
\cfoot{\thepage}

% ============================================================
\begin{document}

\begin{center}
  {\LARGE\bfseries Diario de desarrollo}\\[6pt]
  {\large TFG: Optimizacion SLAM TurtleBot4}\\[4pt]
  {\small Jesús Ruslan Esquinas Herrera \quad·\quad Universidad de Sevilla}\\[2pt]
  {\small Curso 2025--2026}
\end{center}

\vspace{1em}
\hrule
\vspace{1.5em}

% ── Resumen de mapeo (referencia rápida) ─────────────────────
\section*{Mapa de secciones}

\begin{tabular}{>{\ttfamily}p{4.5cm} p{9cm}}
  \toprule
  \normalfont\bfseries Sección & \bfseries Contenido esperado \\
  \midrule
  3.3.1 & Planificación temporal, instalaciones, entorno \\
  3.3.2 & Metodología de desarrollo (iterativa, tutoriales, etc.) \\
  4.2   & Marco teórico: ROS2, Nav2, Create3, algoritmos \\
  4.3   & Requisitos: velocidad máx., frecuencia cmd\_vel, etc. \\
  5.2   & Arquitectura del sistema (simulado y real) \\
  5.2.1 & Diferencias y similitudes simulado vs. real \\
  6.1   & Introducción a la implementación, entorno software \\
  6.2   & Desarrollo de algoritmos (open-loop, closed-loop, Nav2\ldots) \\
  6.3   & Metodología experimental, incidencias y decisiones \\
  7.3   & Resultados obtenidos y tablas comparativas \\
  7.4   & Optimizaciones realizadas \\
  \bottomrule
\end{tabular}

\vspace{1.5em}
\hrule
\vspace{1.5em}

% ============================================================
%  ENTRADAS DEL DIARIO
%  Añade una entrada nueva cada vez que avances algo relevante.
% ============================================================

% ── ENTRADAS YA REGISTRADAS ──────────────────────────────────

\begin{entrada}{2025-XX-XX}{Puesta en marcha del entorno de simulación}
  \problema{Necesidad de reproducir el comportamiento del TurtleBot4 sin disponer del robot físico en todo momento.}
  \solucion{Instalación de \texttt{ros-humble-turtlebot4-simulator} e Ignition Gazebo Fortress. Lanzamiento con:
    \texttt{ros2 launch turtlebot4\_ignition\_bringup turtlebot4\_ignition.launch.py}}
  \secciondestino{6.1 — Introducción a la implementación}
  \nota{La simulación incluye el modelo URDF completo del TurtleBot4 Lite con el Create3 y el sensor LiDAR RPLIDAR A1.}
\end{entrada}


\begin{entrada}{2025-XX-XX}{Teleop con teclado — verificación del entorno}
  \solucion{Verificación del entorno mediante \texttt{teleop\_twist\_keyboard}. El robot responde a comandos en \texttt{/cmd\_vel}.}
  \secciondestino{6.1 — Introducción a la implementación}
  \nota{Primer contacto con el topic \texttt{/cmd\_vel} (\texttt{geometry\_msgs/Twist}), que es la interfaz de control fundamental para todos los algoritmos posteriores.}
\end{entrada}


\begin{entrada}{2025-XX-XX}{Nodo C++ — control open-loop con botones del Create3}
  \problema{Mover el robot una distancia fija ($d = 0{,}5$\,m) mediante los botones físicos del Create3 (\texttt{button\_1}, \texttt{button\_2}).}
  \solucion{Nodo suscrito a \texttt{/interface\_buttons} que publica en \texttt{/cmd\_vel}. Control open-loop: la duración del movimiento se calcula como $t = d / v$.}
  \secciondestino{6.2 — Desarrollo de algoritmos (control básico)}
  \nota{Primera implementación de referencia. El error de posición final de este método será la línea base para comparar con el control closed-loop.}
\end{entrada}


\begin{entrada}{2025-XX-XX}{Problema: el robot para solo tras un único comando de velocidad}
  \problema{El robot avanzaba $\approx 0{,}15$\,m y se detenía sin que se alcanzara la distancia objetivo ni se publicara el mensaje de parada.}
  \causa{El plugin del Create3 en Ignition Gazebo implementa un \textit{safety timeout}: si deja de recibir comandos en \texttt{/cmd\_vel} con regularidad, el robot para automáticamente. Una sola publicación no es suficiente.}
  \solucion{Sustituir la publicación única por un \texttt{wall\_timer} a 10\,Hz que publica continuamente el \texttt{Twist} mientras \texttt{moving\_} sea \texttt{true}. El timer se cancela al alcanzar la distancia o al pulsar Button 2.}

\begin{lstlisting}[language=C++, caption={Timer de publicación continua en /cmd\_vel}]
cmd_vel_timer_ = this->create_wall_timer(
  100ms,
  [this]() {
    if (!moving_) { cmd_vel_timer_->cancel(); return; }
    auto twist = geometry_msgs::msg::Twist();
    twist.linear.x = LINEAR_SPEED;
    cmd_vel_publisher_->publish(twist);
  });
\end{lstlisting}

  \secciondestino{6.3 — Metodología experimental (incidencias)}
  \resultado{El comportamiento del simulador es fiel al robot real: el Create3 físico exige el mismo patrón de publicación continua a $\geq 10$\,Hz.}
\end{entrada}


\begin{entrada}{2025-XX-XX}{Nodo C++ — control closed-loop con odometría}
  \problema{El control open-loop acumula error porque asume velocidad constante, ignorando la rampa de aceleración interna del Create3.}
  \solucion{Suscripción a \texttt{/odom} (\texttt{nav\_msgs/Odometry}). Se guarda la posición de inicio $(x_0, y_0)$ al pulsar Button 1 y se para cuando:
    \[ d = \sqrt{(x - x_0)^2 + (y - y_0)^2} \geq d_{\text{objetivo}} \]}
  \secciondestino{6.2 — Desarrollo de algoritmos (control con realimentación)}
  \resultado{Comparativa open-loop vs. closed-loop: medir el error de posición final $\varepsilon = |d_{\text{real}} - d_{\text{objetivo}}|$ en ambos métodos. Ver §\,7.3.}
\end{entrada}


\begin{entrada}{2025-XX-XX}{Problema: velocidad saturada a $\approx 0{,}31$\,m/s}
  \problema{Al comandar \texttt{LINEAR\_SPEED = 2.0}\,m/s, la odometría reporta $\approx 0{,}31$\,m/s.}
  \causa{El iRobot Create3 impone un límite de velocidad máxima de \textbf{0{,}306\,m/s} por firmware, tanto en robot real como en simulación. Cualquier valor superior se satura.}
  \solucion{Usar \texttt{LINEAR\_SPEED = 0.3}\,m/s como velocidad nominal máxima.}
  \secciondestino{4.3 — Requisitos del sistema (restricciones hardware)}
  \nota{Este límite es un requisito hardware del sistema, no un bug. Debe documentarse en la sección de requisitos.}
\end{entrada}


\begin{entrada}{2025-XX-XX}{Problema: artefactos de compilación con colcon}
  \problema{Tras modificar las constantes \texttt{LINEAR\_SPEED} y \texttt{TARGET\_DIST}, el robot seguía comportándose con los valores antiguos (recorría 5\,m en lugar de 0{,}5\,m).}
  \causa{Los directorios \texttt{build/} e \texttt{install/} cacheaban el binario anterior. \texttt{colcon build} sin limpiar no recompiló el ejecutable.}
  \solucion{Build limpio:
\begin{lstlisting}[language=bash]
rm -rf build/ install/ log/
colcon build --packages-select turtlebot4_cpp_tutorials
source install/setup.bash
\end{lstlisting}}
  \secciondestino{6.3 — Metodología experimental (decisiones técnicas)}
  \nota{Buena práctica: ante comportamientos inesperados, verificar siempre que el binario en ejecución corresponde al código fuente actual.}
\end{entrada}


% ============================================================
%  PLANTILLA PARA NUEVAS ENTRADAS
%  Copia el bloque siguiente cada vez que añadas una entrada.
% ============================================================

% \begin{entrada}{YYYY-MM-DD}{Título descriptivo de la tarea o problema}
%   \problema{Descripción del problema o tarea abordada.}
%   \causa{Causa raíz identificada (si aplica).}
%   \solucion{Qué se hizo para resolverlo.}
%
%   % Incluye código si es relevante:
%   % \begin{lstlisting}[language=C++, caption={Descripción}]
%   % // código aquí
%   % \end{lstlisting}
%
%   \secciondestino{X.Y — Nombre de la sección}
%   \resultado{Dato medible o conclusión extraíble (si aplica).}
%   \nota{Observación adicional o lección aprendida.}
% \end{entrada}

\end{document}

% ============================================================

\documentclass[12pt, a4paper]{article}

\usepackage[utf8]{inputenc}
\usepackage[T1]{fontenc}
\usepackage[spanish]{babel}
\usepackage{geometry}
\geometry{margin=2.5cm}

\usepackage{xcolor}
\usepackage{mdframed}
\usepackage{listings}
\usepackage{hyperref}
\usepackage{parskip}
\usepackage{titlesec}
\usepackage{enumitem}
\usepackage{fancyhdr}
\usepackage{booktabs}
\usepackage{array}

% ── Colores ──────────────────────────────────────────────────
\definecolor{colorProblema}{RGB}{180, 60, 40}
\definecolor{colorSolucion}{RGB}{30, 120, 60}
\definecolor{colorDestino}{RGB}{30, 80, 160}
\definecolor{colorNota}{RGB}{140, 90, 0}
\definecolor{colorFondo}{RGB}{248, 248, 248}
\definecolor{colorBorde}{RGB}{200, 200, 200}
\definecolor{codeGray}{RGB}{245, 245, 245}

% ── Estilo de código ─────────────────────────────────────────
\lstset{
  backgroundcolor=\color{codeGray},
  basicstyle=\ttfamily\small,
  breaklines=true,
  frame=single,
  rulecolor=\color{colorBorde},
  numbers=left,
  numberstyle=\tiny\color{gray},
  tabsize=2,
  showstringspaces=false,
  keywordstyle=\color{colorDestino}\bfseries,
  commentstyle=\color{colorSolucion}\itshape,
  stringstyle=\color{colorProblema},
}

% ── Entorno de entrada de diario ─────────────────────────────
%  \entrada{fecha}{título}
\newenvironment{entrada}[2]{%
  \begin{mdframed}[
    backgroundcolor=colorFondo,
    linecolor=colorBorde,
    linewidth=0.8pt,
    topline=true, bottomline=true,
    leftline=true, rightline=false,
    leftmargin=0pt, rightmargin=0pt,
    innerleftmargin=10pt, innerrightmargin=10pt,
    innertopmargin=8pt, innerbottommargin=8pt,
  ]
  \noindent
  {\small\color{gray}\texttt{#1}} \hfill
  {\small\color{gray}---}\\[2pt]
  {\large\bfseries #2}\\[6pt]
  \hrule height 0.4pt \color{colorBorde}
  \vspace{6pt}
}{%
  \end{mdframed}
  \vspace{8pt}
}

% ── Comandos de etiquetas ─────────────────────────────────────
\newcommand{\problema}[1]{%
  \vspace{4pt}
  \noindent{\color{colorProblema}\bfseries Problema:} #1\par
}

\newcommand{\causa}[1]{%
  \vspace{2pt}
  \noindent{\color{colorProblema}\bfseries Causa:} #1\par
}

\newcommand{\solucion}[1]{%
  \vspace{2pt}
  \noindent{\color{colorSolucion}\bfseries Solución:} #1\par
}

\newcommand{\secciondestino}[1]{%
  \vspace{4pt}
  \noindent{\color{colorDestino}\bfseries $\rightarrow$ Sección en memoria:} \texttt{#1}\par
}

\newcommand{\nota}[1]{%
  \vspace{4pt}
  \noindent{\color{colorNota}\bfseries Nota:} \textit{#1}\par
}

\newcommand{\resultado}[1]{%
  \vspace{4pt}
  \noindent{\color{colorSolucion}\bfseries Resultado medible:} #1\par
}

% ── Cabecera ─────────────────────────────────────────────────
\pagestyle{fancy}
\fancyhf{}
\rhead{\small TFG — Notas de desarrollo}
\lhead{\small Jesús Ruslan}
\cfoot{\thepage}

% ============================================================
\begin{document}

\begin{center}
  {\LARGE\bfseries Diario de desarrollo}\\[6pt]
  {\large TFG: Optimizacion SLAM TurtleBot4}\\[4pt]
  {\small Jesús Ruslan Esquinas Herrera \quad·\quad Universidad de Sevilla}\\[2pt]
  {\small Curso 2025--2026}
\end{center}

\vspace{1em}
\hrule
\vspace{1.5em}

% ── Resumen de mapeo (referencia rápida) ─────────────────────
\section*{Mapa de secciones}

\begin{tabular}{>{\ttfamily}p{4.5cm} p{9cm}}
  \toprule
  \normalfont\bfseries Sección & \bfseries Contenido esperado \\
  \midrule
  3.3.1 & Planificación temporal, instalaciones, entorno \\
  3.3.2 & Metodología de desarrollo (iterativa, tutoriales, etc.) \\
  4.2   & Marco teórico: ROS2, Nav2, Create3, algoritmos \\
  4.3   & Requisitos: velocidad máx., frecuencia cmd\_vel, etc. \\
  5.2   & Arquitectura del sistema (simulado y real) \\
  5.2.1 & Diferencias y similitudes simulado vs. real \\
  6.1   & Introducción a la implementación, entorno software \\
  6.2   & Desarrollo de algoritmos (open-loop, closed-loop, Nav2\ldots) \\
  6.3   & Metodología experimental, incidencias y decisiones \\
  7.3   & Resultados obtenidos y tablas comparativas \\
  7.4   & Optimizaciones realizadas \\
  \bottomrule
\end{tabular}

\vspace{1.5em}
\hrule
\vspace{1.5em}

% ============================================================
%  ENTRADAS DEL DIARIO
%  Añade una entrada nueva cada vez que avances algo relevante.
% ============================================================

% ── ENTRADAS YA REGISTRADAS ──────────────────────────────────

\begin{entrada}{2025-XX-XX}{Puesta en marcha del entorno de simulación}
  \problema{Necesidad de reproducir el comportamiento del TurtleBot4 sin disponer del robot físico en todo momento.}
  \solucion{Instalación de \texttt{ros-humble-turtlebot4-simulator} e Ignition Gazebo Fortress. Lanzamiento con:
    \texttt{ros2 launch turtlebot4\_ignition\_bringup turtlebot4\_ignition.launch.py}}
  \secciondestino{6.1 — Introducción a la implementación}
  \nota{La simulación incluye el modelo URDF completo del TurtleBot4 Lite con el Create3 y el sensor LiDAR RPLIDAR A1.}
\end{entrada}


\begin{entrada}{2025-XX-XX}{Teleop con teclado — verificación del entorno}
  \solucion{Verificación del entorno mediante \texttt{teleop\_twist\_keyboard}. El robot responde a comandos en \texttt{/cmd\_vel}.}
  \secciondestino{6.1 — Introducción a la implementación}
  \nota{Primer contacto con el topic \texttt{/cmd\_vel} (\texttt{geometry\_msgs/Twist}), que es la interfaz de control fundamental para todos los algoritmos posteriores.}
\end{entrada}


\begin{entrada}{2025-XX-XX}{Nodo C++ — control open-loop con botones del Create3}
  \problema{Mover el robot una distancia fija ($d = 0{,}5$\,m) mediante los botones físicos del Create3 (\texttt{button\_1}, \texttt{button\_2}).}
  \solucion{Nodo suscrito a \texttt{/interface\_buttons} que publica en \texttt{/cmd\_vel}. Control open-loop: la duración del movimiento se calcula como $t = d / v$.}
  \secciondestino{6.2 — Desarrollo de algoritmos (control básico)}
  \nota{Primera implementación de referencia. El error de posición final de este método será la línea base para comparar con el control closed-loop.}
\end{entrada}


\begin{entrada}{2025-XX-XX}{Problema: el robot para solo tras un único comando de velocidad}
  \problema{El robot avanzaba $\approx 0{,}15$\,m y se detenía sin que se alcanzara la distancia objetivo ni se publicara el mensaje de parada.}
  \causa{El plugin del Create3 en Ignition Gazebo implementa un \textit{safety timeout}: si deja de recibir comandos en \texttt{/cmd\_vel} con regularidad, el robot para automáticamente. Una sola publicación no es suficiente.}
  \solucion{Sustituir la publicación única por un \texttt{wall\_timer} a 10\,Hz que publica continuamente el \texttt{Twist} mientras \texttt{moving\_} sea \texttt{true}. El timer se cancela al alcanzar la distancia o al pulsar Button 2.}

\begin{lstlisting}[language=C++, caption={Timer de publicación continua en /cmd\_vel}]
cmd_vel_timer_ = this->create_wall_timer(
  100ms,
  [this]() {
    if (!moving_) { cmd_vel_timer_->cancel(); return; }
    auto twist = geometry_msgs::msg::Twist();
    twist.linear.x = LINEAR_SPEED;
    cmd_vel_publisher_->publish(twist);
  });
\end{lstlisting}

  \secciondestino{6.3 — Metodología experimental (incidencias)}
  \resultado{El comportamiento del simulador es fiel al robot real: el Create3 físico exige el mismo patrón de publicación continua a $\geq 10$\,Hz.}
\end{entrada}


\begin{entrada}{2025-XX-XX}{Nodo C++ — control closed-loop con odometría}
  \problema{El control open-loop acumula error porque asume velocidad constante, ignorando la rampa de aceleración interna del Create3.}
  \solucion{Suscripción a \texttt{/odom} (\texttt{nav\_msgs/Odometry}). Se guarda la posición de inicio $(x_0, y_0)$ al pulsar Button 1 y se para cuando:
    \[ d = \sqrt{(x - x_0)^2 + (y - y_0)^2} \geq d_{\text{objetivo}} \]}
  \secciondestino{6.2 — Desarrollo de algoritmos (control con realimentación)}
  \resultado{Comparativa open-loop vs. closed-loop: medir el error de posición final $\varepsilon = |d_{\text{real}} - d_{\text{objetivo}}|$ en ambos métodos. Ver §\,7.3.}
\end{entrada}


\begin{entrada}{2025-XX-XX}{Problema: velocidad saturada a $\approx 0{,}31$\,m/s}
  \problema{Al comandar \texttt{LINEAR\_SPEED = 2.0}\,m/s, la odometría reporta $\approx 0{,}31$\,m/s.}
  \causa{El iRobot Create3 impone un límite de velocidad máxima de \textbf{0{,}306\,m/s} por firmware, tanto en robot real como en simulación. Cualquier valor superior se satura.}
  \solucion{Usar \texttt{LINEAR\_SPEED = 0.3}\,m/s como velocidad nominal máxima.}
  \secciondestino{4.3 — Requisitos del sistema (restricciones hardware)}
  \nota{Este límite es un requisito hardware del sistema, no un bug. Debe documentarse en la sección de requisitos.}
\end{entrada}


\begin{entrada}{2025-XX-XX}{Problema: artefactos de compilación con colcon}
  \problema{Tras modificar las constantes \texttt{LINEAR\_SPEED} y \texttt{TARGET\_DIST}, el robot seguía comportándose con los valores antiguos (recorría 5\,m en lugar de 0{,}5\,m).}
  \causa{Los directorios \texttt{build/} e \texttt{install/} cacheaban el binario anterior. \texttt{colcon build} sin limpiar no recompiló el ejecutable.}
  \solucion{Build limpio:
\begin{lstlisting}[language=bash]
rm -rf build/ install/ log/
colcon build --packages-select turtlebot4_cpp_tutorials
source install/setup.bash
\end{lstlisting}}
  \secciondestino{6.3 — Metodología experimental (decisiones técnicas)}
  \nota{Buena práctica: ante comportamientos inesperados, verificar siempre que el binario en ejecución corresponde al código fuente actual.}
\end{entrada}


% ============================================================
%  PLANTILLA PARA NUEVAS ENTRADAS
%  Copia el bloque siguiente cada vez que añadas una entrada.
% ============================================================

% \begin{entrada}{YYYY-MM-DD}{Título descriptivo de la tarea o problema}
%   \problema{Descripción del problema o tarea abordada.}
%   \causa{Causa raíz identificada (si aplica).}
%   \solucion{Qué se hizo para resolverlo.}
%
%   % Incluye código si es relevante:
%   % \begin{lstlisting}[language=C++, caption={Descripción}]
%   % // código aquí
%   % \end{lstlisting}
%
%   \secciondestino{X.Y — Nombre de la sección}
%   \resultado{Dato medible o conclusión extraíble (si aplica).}
%   \nota{Observación adicional o lección aprendida.}
% \end{entrada}

\end{document}

% ============================================================

\documentclass[12pt, a4paper]{article}

\usepackage[utf8]{inputenc}
\usepackage[T1]{fontenc}
\usepackage[spanish]{babel}
\usepackage{geometry}
\geometry{margin=2.5cm}

\usepackage{xcolor}
\usepackage{mdframed}
\usepackage{listings}
\usepackage{hyperref}
\usepackage{parskip}
\usepackage{titlesec}
\usepackage{enumitem}
\usepackage{fancyhdr}
\usepackage{booktabs}
\usepackage{array}

% ── Colores ──────────────────────────────────────────────────
\definecolor{colorProblema}{RGB}{180, 60, 40}
\definecolor{colorSolucion}{RGB}{30, 120, 60}
\definecolor{colorDestino}{RGB}{30, 80, 160}
\definecolor{colorNota}{RGB}{140, 90, 0}
\definecolor{colorFondo}{RGB}{248, 248, 248}
\definecolor{colorBorde}{RGB}{200, 200, 200}
\definecolor{codeGray}{RGB}{245, 245, 245}

% ── Estilo de código ─────────────────────────────────────────
\lstset{
  backgroundcolor=\color{codeGray},
  basicstyle=\ttfamily\small,
  breaklines=true,
  frame=single,
  rulecolor=\color{colorBorde},
  numbers=left,
  numberstyle=\tiny\color{gray},
  tabsize=2,
  showstringspaces=false,
  keywordstyle=\color{colorDestino}\bfseries,
  commentstyle=\color{colorSolucion}\itshape,
  stringstyle=\color{colorProblema},
}

% ── Entorno de entrada de diario ─────────────────────────────
%  \entrada{fecha}{título}
\newenvironment{entrada}[2]{%
  \begin{mdframed}[
    backgroundcolor=colorFondo,
    linecolor=colorBorde,
    linewidth=0.8pt,
    topline=true, bottomline=true,
    leftline=true, rightline=false,
    leftmargin=0pt, rightmargin=0pt,
    innerleftmargin=10pt, innerrightmargin=10pt,
    innertopmargin=8pt, innerbottommargin=8pt,
  ]
  \noindent
  {\small\color{gray}\texttt{#1}} \hfill
  {\small\color{gray}---}\\[2pt]
  {\large\bfseries #2}\\[6pt]
  \hrule height 0.4pt \color{colorBorde}
  \vspace{6pt}
}{%
  \end{mdframed}
  \vspace{8pt}
}

% ── Comandos de etiquetas ─────────────────────────────────────
\newcommand{\problema}[1]{%
  \vspace{4pt}
  \noindent{\color{colorProblema}\bfseries Problema:} #1\par
}

\newcommand{\causa}[1]{%
  \vspace{2pt}
  \noindent{\color{colorProblema}\bfseries Causa:} #1\par
}

\newcommand{\solucion}[1]{%
  \vspace{2pt}
  \noindent{\color{colorSolucion}\bfseries Solución:} #1\par
}

\newcommand{\secciondestino}[1]{%
  \vspace{4pt}
  \noindent{\color{colorDestino}\bfseries $\rightarrow$ Sección en memoria:} \texttt{#1}\par
}

\newcommand{\nota}[1]{%
  \vspace{4pt}
  \noindent{\color{colorNota}\bfseries Nota:} \textit{#1}\par
}

\newcommand{\resultado}[1]{%
  \vspace{4pt}
  \noindent{\color{colorSolucion}\bfseries Resultado medible:} #1\par
}

% ── Cabecera ─────────────────────────────────────────────────
\pagestyle{fancy}
\fancyhf{}
\rhead{\small TFG — Notas de desarrollo}
\lhead{\small Jesús Ruslan}
\cfoot{\thepage}

% ============================================================
\begin{document}

\begin{center}
  {\LARGE\bfseries Diario de desarrollo}\\[6pt]
  {\large TFG: Optimizacion SLAM TurtleBot4}\\[4pt]
  {\small Jesús Ruslan Esquinas Herrera \quad·\quad Universidad de Sevilla}\\[2pt]
  {\small Curso 2025--2026}
\end{center}

\vspace{1em}
\hrule
\vspace{1.5em}

% ── Resumen de mapeo (referencia rápida) ─────────────────────
\section*{Mapa de secciones}

\begin{tabular}{>{\ttfamily}p{4.5cm} p{9cm}}
  \toprule
  \normalfont\bfseries Sección & \bfseries Contenido esperado \\
  \midrule
  3.3.1 & Planificación temporal, instalaciones, entorno \\
  3.3.2 & Metodología de desarrollo (iterativa, tutoriales, etc.) \\
  4.2   & Marco teórico: ROS2, Nav2, Create3, algoritmos \\
  4.3   & Requisitos: velocidad máx., frecuencia cmd\_vel, etc. \\
  5.2   & Arquitectura del sistema (simulado y real) \\
  5.2.1 & Diferencias y similitudes simulado vs. real \\
  6.1   & Introducción a la implementación, entorno software \\
  6.2   & Desarrollo de algoritmos (open-loop, closed-loop, Nav2\ldots) \\
  6.3   & Metodología experimental, incidencias y decisiones \\
  7.3   & Resultados obtenidos y tablas comparativas \\
  7.4   & Optimizaciones realizadas \\
  \bottomrule
\end{tabular}

\vspace{1.5em}
\hrule
\vspace{1.5em}

% ============================================================
%  ENTRADAS DEL DIARIO
%  Añade una entrada nueva cada vez que avances algo relevante.
% ============================================================

% ── ENTRADAS YA REGISTRADAS ──────────────────────────────────

\begin{entrada}{2025-XX-XX}{Puesta en marcha del entorno de simulación}
  \problema{Necesidad de reproducir el comportamiento del TurtleBot4 sin disponer del robot físico en todo momento.}
  \solucion{Instalación de \texttt{ros-humble-turtlebot4-simulator} e Ignition Gazebo Fortress. Lanzamiento con:
    \texttt{ros2 launch turtlebot4\_ignition\_bringup turtlebot4\_ignition.launch.py}}
  \secciondestino{6.1 — Introducción a la implementación}
  \nota{La simulación incluye el modelo URDF completo del TurtleBot4 Lite con el Create3 y el sensor LiDAR RPLIDAR A1.}
\end{entrada}


\begin{entrada}{2025-XX-XX}{Teleop con teclado — verificación del entorno}
  \solucion{Verificación del entorno mediante \texttt{teleop\_twist\_keyboard}. El robot responde a comandos en \texttt{/cmd\_vel}.}
  \secciondestino{6.1 — Introducción a la implementación}
  \nota{Primer contacto con el topic \texttt{/cmd\_vel} (\texttt{geometry\_msgs/Twist}), que es la interfaz de control fundamental para todos los algoritmos posteriores.}
\end{entrada}


\begin{entrada}{2025-XX-XX}{Nodo C++ — control open-loop con botones del Create3}
  \problema{Mover el robot una distancia fija ($d = 0{,}5$\,m) mediante los botones físicos del Create3 (\texttt{button\_1}, \texttt{button\_2}).}
  \solucion{Nodo suscrito a \texttt{/interface\_buttons} que publica en \texttt{/cmd\_vel}. Control open-loop: la duración del movimiento se calcula como $t = d / v$.}
  \secciondestino{6.2 — Desarrollo de algoritmos (control básico)}
  \nota{Primera implementación de referencia. El error de posición final de este método será la línea base para comparar con el control closed-loop.}
\end{entrada}


\begin{entrada}{2025-XX-XX}{Problema: el robot para solo tras un único comando de velocidad}
  \problema{El robot avanzaba $\approx 0{,}15$\,m y se detenía sin que se alcanzara la distancia objetivo ni se publicara el mensaje de parada.}
  \causa{El plugin del Create3 en Ignition Gazebo implementa un \textit{safety timeout}: si deja de recibir comandos en \texttt{/cmd\_vel} con regularidad, el robot para automáticamente. Una sola publicación no es suficiente.}
  \solucion{Sustituir la publicación única por un \texttt{wall\_timer} a 10\,Hz que publica continuamente el \texttt{Twist} mientras \texttt{moving\_} sea \texttt{true}. El timer se cancela al alcanzar la distancia o al pulsar Button 2.}

\begin{lstlisting}[language=C++, caption={Timer de publicación continua en /cmd\_vel}]
cmd_vel_timer_ = this->create_wall_timer(
  100ms,
  [this]() {
    if (!moving_) { cmd_vel_timer_->cancel(); return; }
    auto twist = geometry_msgs::msg::Twist();
    twist.linear.x = LINEAR_SPEED;
    cmd_vel_publisher_->publish(twist);
  });
\end{lstlisting}

  \secciondestino{6.3 — Metodología experimental (incidencias)}
  \resultado{El comportamiento del simulador es fiel al robot real: el Create3 físico exige el mismo patrón de publicación continua a $\geq 10$\,Hz.}
\end{entrada}


\begin{entrada}{2025-XX-XX}{Nodo C++ — control closed-loop con odometría}
  \problema{El control open-loop acumula error porque asume velocidad constante, ignorando la rampa de aceleración interna del Create3.}
  \solucion{Suscripción a \texttt{/odom} (\texttt{nav\_msgs/Odometry}). Se guarda la posición de inicio $(x_0, y_0)$ al pulsar Button 1 y se para cuando:
    \[ d = \sqrt{(x - x_0)^2 + (y - y_0)^2} \geq d_{\text{objetivo}} \]}
  \secciondestino{6.2 — Desarrollo de algoritmos (control con realimentación)}
  \resultado{Comparativa open-loop vs. closed-loop: medir el error de posición final $\varepsilon = |d_{\text{real}} - d_{\text{objetivo}}|$ en ambos métodos. Ver §\,7.3.}
\end{entrada}


\begin{entrada}{2025-XX-XX}{Problema: velocidad saturada a $\approx 0{,}31$\,m/s}
  \problema{Al comandar \texttt{LINEAR\_SPEED = 2.0}\,m/s, la odometría reporta $\approx 0{,}31$\,m/s.}
  \causa{El iRobot Create3 impone un límite de velocidad máxima de \textbf{0{,}306\,m/s} por firmware, tanto en robot real como en simulación. Cualquier valor superior se satura.}
  \solucion{Usar \texttt{LINEAR\_SPEED = 0.3}\,m/s como velocidad nominal máxima.}
  \secciondestino{4.3 — Requisitos del sistema (restricciones hardware)}
  \nota{Este límite es un requisito hardware del sistema, no un bug. Debe documentarse en la sección de requisitos.}
\end{entrada}


\begin{entrada}{2025-XX-XX}{Problema: artefactos de compilación con colcon}
  \problema{Tras modificar las constantes \texttt{LINEAR\_SPEED} y \texttt{TARGET\_DIST}, el robot seguía comportándose con los valores antiguos (recorría 5\,m en lugar de 0{,}5\,m).}
  \causa{Los directorios \texttt{build/} e \texttt{install/} cacheaban el binario anterior. \texttt{colcon build} sin limpiar no recompiló el ejecutable.}
  \solucion{Build limpio:
\begin{lstlisting}[language=bash]
rm -rf build/ install/ log/
colcon build --packages-select turtlebot4_cpp_tutorials
source install/setup.bash
\end{lstlisting}}
  \secciondestino{6.3 — Metodología experimental (decisiones técnicas)}
  \nota{Buena práctica: ante comportamientos inesperados, verificar siempre que el binario en ejecución corresponde al código fuente actual.}
\end{entrada}


% ============================================================
%  PLANTILLA PARA NUEVAS ENTRADAS
%  Copia el bloque siguiente cada vez que añadas una entrada.
% ============================================================

% \begin{entrada}{YYYY-MM-DD}{Título descriptivo de la tarea o problema}
%   \problema{Descripción del problema o tarea abordada.}
%   \causa{Causa raíz identificada (si aplica).}
%   \solucion{Qué se hizo para resolverlo.}
%
%   % Incluye código si es relevante:
%   % \begin{lstlisting}[language=C++, caption={Descripción}]
%   % // código aquí
%   % \end{lstlisting}
%
%   \secciondestino{X.Y — Nombre de la sección}
%   \resultado{Dato medible o conclusión extraíble (si aplica).}
%   \nota{Observación adicional o lección aprendida.}
% \end{entrada}

\end{document}

% ============================================================

\documentclass[12pt, a4paper]{article}

\usepackage[utf8]{inputenc}
\usepackage[T1]{fontenc}
\usepackage[spanish]{babel}
\usepackage{geometry}
\geometry{margin=2.5cm}

\usepackage{xcolor}
\usepackage{mdframed}
\usepackage{listings}
\usepackage{hyperref}
\usepackage{parskip}
\usepackage{titlesec}
\usepackage{enumitem}
\usepackage{fancyhdr}
\usepackage{booktabs}
\usepackage{array}

% ── Colores ──────────────────────────────────────────────────
\definecolor{colorProblema}{RGB}{180, 60, 40}
\definecolor{colorSolucion}{RGB}{30, 120, 60}
\definecolor{colorDestino}{RGB}{30, 80, 160}
\definecolor{colorNota}{RGB}{140, 90, 0}
\definecolor{colorFondo}{RGB}{248, 248, 248}
\definecolor{colorBorde}{RGB}{200, 200, 200}
\definecolor{codeGray}{RGB}{245, 245, 245}

% ── Estilo de código ─────────────────────────────────────────
\lstset{
  backgroundcolor=\color{codeGray},
  basicstyle=\ttfamily\small,
  breaklines=true,
  frame=single,
  rulecolor=\color{colorBorde},
  numbers=left,
  numberstyle=\tiny\color{gray},
  tabsize=2,
  showstringspaces=false,
  keywordstyle=\color{colorDestino}\bfseries,
  commentstyle=\color{colorSolucion}\itshape,
  stringstyle=\color{colorProblema},
}

% ── Entorno de entrada de diario ─────────────────────────────
%  \entrada{fecha}{título}
\newenvironment{entrada}[2]{%
  \begin{mdframed}[
    backgroundcolor=colorFondo,
    linecolor=colorBorde,
    linewidth=0.8pt,
    topline=true, bottomline=true,
    leftline=true, rightline=false,
    leftmargin=0pt, rightmargin=0pt,
    innerleftmargin=10pt, innerrightmargin=10pt,
    innertopmargin=8pt, innerbottommargin=8pt,
  ]
  \noindent
  {\small\color{gray}\texttt{#1}} \hfill
  {\small\color{gray}---}\\[2pt]
  {\large\bfseries #2}\\[6pt]
  \hrule height 0.4pt \color{colorBorde}
  \vspace{6pt}
}{%
  \end{mdframed}
  \vspace{8pt}
}

% ── Comandos de etiquetas ─────────────────────────────────────
\newcommand{\problema}[1]{%
  \vspace{4pt}
  \noindent{\color{colorProblema}\bfseries Problema:} #1\par
}

\newcommand{\causa}[1]{%
  \vspace{2pt}
  \noindent{\color{colorProblema}\bfseries Causa:} #1\par
}

\newcommand{\solucion}[1]{%
  \vspace{2pt}
  \noindent{\color{colorSolucion}\bfseries Solución:} #1\par
}

\newcommand{\secciondestino}[1]{%
  \vspace{4pt}
  \noindent{\color{colorDestino}\bfseries $\rightarrow$ Sección en memoria:} \texttt{#1}\par
}

\newcommand{\nota}[1]{%
  \vspace{4pt}
  \noindent{\color{colorNota}\bfseries Nota:} \textit{#1}\par
}

\newcommand{\resultado}[1]{%
  \vspace{4pt}
  \noindent{\color{colorSolucion}\bfseries Resultado medible:} #1\par
}

% ── Cabecera ─────────────────────────────────────────────────
\pagestyle{fancy}
\fancyhf{}
\rhead{\small TFG — Notas de desarrollo}
\lhead{\small Jesús Ruslan}
\cfoot{\thepage}

% ============================================================
\begin{document}

\begin{center}
  {\LARGE\bfseries Diario de desarrollo}\\[6pt]
  {\large TFG: Optimizacion SLAM TurtleBot4}\\[4pt]
  {\small Jesús Ruslan Esquinas Herrera \quad·\quad Universidad de Sevilla}\\[2pt]
  {\small Curso 2025--2026}
\end{center}

\vspace{1em}
\hrule
\vspace{1.5em}

% ── Resumen de mapeo (referencia rápida) ─────────────────────
\section*{Mapa de secciones}

\begin{tabular}{>{\ttfamily}p{4.5cm} p{9cm}}
  \toprule
  \normalfont\bfseries Sección & \bfseries Contenido esperado \\
  \midrule
  3.3.1 & Planificación temporal, instalaciones, entorno \\
  3.3.2 & Metodología de desarrollo (iterativa, tutoriales, etc.) \\
  4.2   & Marco teórico: ROS2, Nav2, Create3, algoritmos \\
  4.3   & Requisitos: velocidad máx., frecuencia cmd\_vel, etc. \\
  5.2   & Arquitectura del sistema (simulado y real) \\
  5.2.1 & Diferencias y similitudes simulado vs. real \\
  6.1   & Introducción a la implementación, entorno software \\
  6.2   & Desarrollo de algoritmos (open-loop, closed-loop, Nav2\ldots) \\
  6.3   & Metodología experimental, incidencias y decisiones \\
  7.3   & Resultados obtenidos y tablas comparativas \\
  7.4   & Optimizaciones realizadas \\
  \bottomrule
\end{tabular}

\vspace{1.5em}
\hrule
\vspace{1.5em}

% ============================================================
%  ENTRADAS DEL DIARIO
%  Añade una entrada nueva cada vez que avances algo relevante.
% ============================================================

% ── ENTRADAS YA REGISTRADAS ──────────────────────────────────

\begin{entrada}{2025-XX-XX}{Puesta en marcha del entorno de simulación}
  \problema{Necesidad de reproducir el comportamiento del TurtleBot4 sin disponer del robot físico en todo momento.}
  \solucion{Instalación de \texttt{ros-humble-turtlebot4-simulator} e Ignition Gazebo Fortress. Lanzamiento con:
    \texttt{ros2 launch turtlebot4\_ignition\_bringup turtlebot4\_ignition.launch.py}}
  \secciondestino{6.1 — Introducción a la implementación}
  \nota{La simulación incluye el modelo URDF completo del TurtleBot4 Lite con el Create3 y el sensor LiDAR RPLIDAR A1.}
\end{entrada}


\begin{entrada}{2025-XX-XX}{Teleop con teclado — verificación del entorno}
  \solucion{Verificación del entorno mediante \texttt{teleop\_twist\_keyboard}. El robot responde a comandos en \texttt{/cmd\_vel}.}
  \secciondestino{6.1 — Introducción a la implementación}
  \nota{Primer contacto con el topic \texttt{/cmd\_vel} (\texttt{geometry\_msgs/Twist}), que es la interfaz de control fundamental para todos los algoritmos posteriores.}
\end{entrada}


\begin{entrada}{2025-XX-XX}{Nodo C++ — control open-loop con botones del Create3}
  \problema{Mover el robot una distancia fija ($d = 0{,}5$\,m) mediante los botones físicos del Create3 (\texttt{button\_1}, \texttt{button\_2}).}
  \solucion{Nodo suscrito a \texttt{/interface\_buttons} que publica en \texttt{/cmd\_vel}. Control open-loop: la duración del movimiento se calcula como $t = d / v$.}
  \secciondestino{6.2 — Desarrollo de algoritmos (control básico)}
  \nota{Primera implementación de referencia. El error de posición final de este método será la línea base para comparar con el control closed-loop.}
\end{entrada}


\begin{entrada}{2025-XX-XX}{Problema: el robot para solo tras un único comando de velocidad}
  \problema{El robot avanzaba $\approx 0{,}15$\,m y se detenía sin que se alcanzara la distancia objetivo ni se publicara el mensaje de parada.}
  \causa{El plugin del Create3 en Ignition Gazebo implementa un \textit{safety timeout}: si deja de recibir comandos en \texttt{/cmd\_vel} con regularidad, el robot para automáticamente. Una sola publicación no es suficiente.}
  \solucion{Sustituir la publicación única por un \texttt{wall\_timer} a 10\,Hz que publica continuamente el \texttt{Twist} mientras \texttt{moving\_} sea \texttt{true}. El timer se cancela al alcanzar la distancia o al pulsar Button 2.}

\begin{lstlisting}[language=C++, caption={Timer de publicación continua en /cmd\_vel}]
cmd_vel_timer_ = this->create_wall_timer(
  100ms,
  [this]() {
    if (!moving_) { cmd_vel_timer_->cancel(); return; }
    auto twist = geometry_msgs::msg::Twist();
    twist.linear.x = LINEAR_SPEED;
    cmd_vel_publisher_->publish(twist);
  });
\end{lstlisting}

  \secciondestino{6.3 — Metodología experimental (incidencias)}
  \resultado{El comportamiento del simulador es fiel al robot real: el Create3 físico exige el mismo patrón de publicación continua a $\geq 10$\,Hz.}
\end{entrada}


\begin{entrada}{2025-XX-XX}{Nodo C++ — control closed-loop con odometría}
  \problema{El control open-loop acumula error porque asume velocidad constante, ignorando la rampa de aceleración interna del Create3.}
  \solucion{Suscripción a \texttt{/odom} (\texttt{nav\_msgs/Odometry}). Se guarda la posición de inicio $(x_0, y_0)$ al pulsar Button 1 y se para cuando:
    \[ d = \sqrt{(x - x_0)^2 + (y - y_0)^2} \geq d_{\text{objetivo}} \]}
  \secciondestino{6.2 — Desarrollo de algoritmos (control con realimentación)}
  \resultado{Comparativa open-loop vs. closed-loop: medir el error de posición final $\varepsilon = |d_{\text{real}} - d_{\text{objetivo}}|$ en ambos métodos. Ver §\,7.3.}
\end{entrada}


\begin{entrada}{2025-XX-XX}{Problema: velocidad saturada a $\approx 0{,}31$\,m/s}
  \problema{Al comandar \texttt{LINEAR\_SPEED = 2.0}\,m/s, la odometría reporta $\approx 0{,}31$\,m/s.}
  \causa{El iRobot Create3 impone un límite de velocidad máxima de \textbf{0{,}306\,m/s} por firmware, tanto en robot real como en simulación. Cualquier valor superior se satura.}
  \solucion{Usar \texttt{LINEAR\_SPEED = 0.3}\,m/s como velocidad nominal máxima.}
  \secciondestino{4.3 — Requisitos del sistema (restricciones hardware)}
  \nota{Este límite es un requisito hardware del sistema, no un bug. Debe documentarse en la sección de requisitos.}
\end{entrada}


\begin{entrada}{2025-XX-XX}{Problema: artefactos de compilación con colcon}
  \problema{Tras modificar las constantes \texttt{LINEAR\_SPEED} y \texttt{TARGET\_DIST}, el robot seguía comportándose con los valores antiguos (recorría 5\,m en lugar de 0{,}5\,m).}
  \causa{Los directorios \texttt{build/} e \texttt{install/} cacheaban el binario anterior. \texttt{colcon build} sin limpiar no recompiló el ejecutable.}
  \solucion{Build limpio:
\begin{lstlisting}[language=bash]
rm -rf build/ install/ log/
colcon build --packages-select turtlebot4_cpp_tutorials
source install/setup.bash
\end{lstlisting}}
  \secciondestino{6.3 — Metodología experimental (decisiones técnicas)}
  \nota{Buena práctica: ante comportamientos inesperados, verificar siempre que el binario en ejecución corresponde al código fuente actual.}
\end{entrada}


% ============================================================
%  PLANTILLA PARA NUEVAS ENTRADAS
%  Copia el bloque siguiente cada vez que añadas una entrada.
% ============================================================

% \begin{entrada}{YYYY-MM-DD}{Título descriptivo de la tarea o problema}
%   \problema{Descripción del problema o tarea abordada.}
%   \causa{Causa raíz identificada (si aplica).}
%   \solucion{Qué se hizo para resolverlo.}
%
%   % Incluye código si es relevante:
%   % \begin{lstlisting}[language=C++, caption={Descripción}]
%   % // código aquí
%   % \end{lstlisting}
%
%   \secciondestino{X.Y — Nombre de la sección}
%   \resultado{Dato medible o conclusión extraíble (si aplica).}
%   \nota{Observación adicional o lección aprendida.}
% \end{entrada}

\end{document}
